% !PRINTED 512
\weberpage{512}

1. Wenn die conjugirten Gruppen $\pi^{-1}Q\pi$ theilerfremd sind,
so ist die gleichzeitige Adjunction s\"ammtlicher Wurzeln der
Resolvente $\varphi(t) = 0$ (nach Satz 4., \S.\,155) gleichbedeutend mit
der Adjunction einer zur Einheitsgruppe geh\"origen Function,
und die L\"osung der gegebenen Gleichung ist auf die vollst\"andige
L\"osung der Resolvente zur\"uckgef\"uhrt. Es ist
\[ N = \Omega(\psi, \psi_1 \ldots \psi_{j-1}), \]
und eine \gesperrt{Galois}'sche Resolvente der Gleichung $\varphi(t) = 0$ ist
zugleich eine \gesperrt{Galois}'sche Resolvente der urspr\"unglichen Gleichung.
In diesem Falle nennen wir $\varphi(t) = 0$ eine \gesperrt{Totalresolvente}
der gegebenen Gleichung.

2. Haben die conjugirten Gruppen $\pi^{-1}Q\pi$ einen von der
Einheitsgruppe verschiedenen Theiler $R$ vom Grade $r$, der dann
ein Normaltheiler von $P$ ist, so ist die gleichzeitige Adjunction
von s\"ammtlichen zu $\psi$ conjugirten Functionen gleichbedeutend
mit der Adjunction einer zu $R$ geh\"origen Gr\"osse. Die \gesperrt{Galois}'sche
Resolvente der gegebenen Gleichung ist durch diese Adjunction
noch nicht vollst\"andig gel\"ost, sondern sie ist nur in Factoren
vom Grade $r$ zerlegt. In diesem Falle heisst $\varphi(t) = 0$ eine
\gesperrt{Partialresolvente}.

Ist $P$ eine \gesperrt{einfache Gruppe}, so existiren \gesperrt{nur} Totalresolventen,
w\"ahrend, wenn $P$ Normaltheiler hat, zu jedem
solchen Normaltheiler eine Partialresolvente gefunden werden
kann.

Derselbe Unterschied tritt auch hervor, wenn wir die
\gesperrt{Galois}'sche Gruppe der Resolvente $\varphi(t)$ untersuchen. Diese
Gruppe besteht aus allen Vertauschungen, die in der Reihe der
Gr\"ossen
\weq{1}{\psi, \psi_1, \psi_2 \ldots \psi_{j-1}}
durch Anwendung der s\"ammtlichen Operationen $\pi$ von $P$ hervorgerufen werden; denn jede Gleichung in $\Omega$ zwischen den Gr\"ossen (1)
bleibt richtig, wenn eine solche Permutation vorgenommen wird,
da man die Operationen $\pi$ auf eine solche Gleichung anwenden
kann; und wenn eine Function in $\Omega$ der Gr\"ossen (1) alle diese
Permutationen gestattet, so gestattet sie auch alle Permutationen $\pi$
und ist also gleich einer Gr\"osse in $\Omega$.

Es ist nun der Grad dieser Gruppe zu bestimmen. Unter
den Operationen von $P$ werden die und nur die unter den
Gr\"ossen (1) keine Ver\"anderung hervorrufen, die gleichzeitig in
%CONT>>
